\chapter*{Resume}
\addcontentsline{toc}{chapter}{Resume}
PerTrack est une application web de gestion de la performance et d'evaluation RH. Le systeme implemente permet de centraliser les cycles annuels, les objectifs, les KPI, les equipes, les taches, les points de suivi, les evaluations finales, la validation RH, les plans d'amelioration, les decisions de bonus ou de penalite, les recommandations de carriere, les tableaux de bord, les journaux d'audit, les notifications et certaines fonctions d'assistance par intelligence artificielle.

Le projet repose sur une application frontend React avec Vite, une API backend Node.js et Express, une base MongoDB modelisee avec Mongoose, et un service Python Flask pour les predictions de performance. Le depot contient aussi des artefacts Docker, Kubernetes, Kustomize, Argo CD et GitHub Actions. Ce rapport presente l'analyse fonctionnelle, technique, architecturale, securite, test, IA et DevOps du projet en s'appuyant uniquement sur les elements verifies dans le depot.

\textbf{Mots-cles:} gestion de la performance, evaluation RH, MERN, MongoDB, React, Express, intelligence artificielle, XGBoost, Kubernetes, GitOps.
\clearpage
